\subsection{最初の共有係数表示七族の独立排除}
\label{subsec:cubic-as4-first-shared}
\cite{AS4lowpinned}のFN2--FN3から、次の具体表示だけを参照する。
前節と同じ $U,K,A$ を用い、
\[
\begin{gathered}
 x_1=1,\quad x_2=e,\quad x_5=f,\quad x_9=r,\quad
 x_{11}=-c,\quad x_{15}=c,\\
 x_{19}=\lambda/16-c^2,\qquad x_{20}=-e/16,
 \qquad\lambda\ne0,-1
\end{gathered}
\]
として他の $x_j$ を零にした六関係の代数を $B(\lambda;c,e,f,r)$ と書く。
この記号を全パラメータでAS4と仮定するわけではない。
以下は指定した表示の独立なH3障害であり、全射影型の被覆や
AS4存在十分性・全同型は新たに輸入しない。

\begin{proposition}\label{prop:cubic-fn-lm}
$h\ne\pm1$, $f$ 任意とし、$B=B(1;h/4,0,f,0)$ とする。
この代数は整域で $Z(B)_3=0$ であり、H3包絡にならない。
特にFN3の全 $L_c$ ($f=0$) と $M_c$ ($f=1$) を排除する。
\end{proposition}
\begin{proof}
\[
 q=\frac{h-1}{h+1},\quad
 \alpha=\frac{h+1}{4},\quad\beta=\frac4{h-1},\quad
 \gamma=-\frac{h-1}{4},\quad\delta=-\frac4{h+1},\quad
 t=-\frac{4f}{h-1}
\]
と置くと、六関係は
\[
 da=-ad,\quad ba=\alpha db,\quad bd=\beta ab,\quad
 ca=\gamma dc,\quad cd=\delta ac,\quad cb=q^{-1}bc+tad.
\]
順序 $a<d<b<c$ の四つの三次重なりは全て解消し、
PBW基底 $a^id^lb^jc^k$ を得る。
実際、量子平面 $k\langle a,d\rangle/(da+ad)$ にまず
$(a,d)\mapsto(\alpha d,\beta a)$ のautomorphismで $b$ を付加し、
次に $(a,d,b)\mapsto(\gamma d,\delta a,q^{-1}b)$ と
$\partial(b)=tad$, $\partial(a)=\partial(d)=0$ で $c$ を付加したiterated Ore表示である。
従って整域である。六関係と元の表示の対応も全係数で可逆に検算した。

正規二次元 $N=ad$ のnormalizerは、順序 $(a,d,b,c)$ で
\[
 \nu_N=\operatorname{diag}(-1,-1,-q,-q^{-1}).
\]
中心三次 $F$ は $\nu_N(F)=F$ なので、その各PBW単項式は
\[
 i+l+j+k=3,\qquad q^{j-k}=-1
\]
を満たす。$q\ne0,1$ で $|j-k|\le3$ だから、残るのは
\[
 h=0,\qquad h^2=-1,\qquad h^2=-3
\]
だけである。具体的に
\[
 q+1=\frac{2h}{h+1},\quad q^2+1=\frac{2(h^2+1)}{(h+1)^2},\quad
 q^3+1=\frac{2h(h^2+3)}{(h+1)^3}.
\]
$h=0$ では、非零二次元 $bc-2f ad$ のnormalizerが $-I$ である。
従ってこの場合は全奇数中心が零となる。

$h^2=-1$ では $q$ は四乗根であり、候補は
\[
 F=(u a+v d)b^2+(w a+z d)c^2
\]
に限る。$b^2a=q^{-1}ab^2$, $c^2a=qac^2$ と $da=-ad$ を用いると、
$[F,a]=0$ の四係数はそれぞれ $q^{\pm1}-1$ または $q^{\pm1}+1$ の非零倍である。
よって $u=v=w=z=0$。
$h^2=-3$ では $q$ は原始六乗根で、候補は $F=u b^3+v c^3$。
\[
 b^3a=q^{-1}\alpha db^3,\qquad c^3a=q\gamma dc^3
\]
なので、PBW独立性と $[F,a]=0$ から再び $u=v=0$ となる。
これらの二つの共鳴での交換子の単射性は、
$\mathbb Q(i)$ と $\mathbb Q(\sqrt{-3})$ 上でも独立に確認した。

最後に $R A R^{\mathsf T}$ は全 $h,f$ で非零反巡回最高tensorとなる。
従ってAS4実現では $\mu=I$、H3候補は中心三次でなければならない。
その消滅が主張を与える。
\end{proof}

\begin{proposition}\label{prop:cubic-fn-remaining}
FN3の次の指定表示も全てH3包絡から排除される。
\[
\begin{array}{c|c|c}
 \text{表示}&B(\lambda;c,e,f,r)&\text{範囲}\\\hline
 F^\varepsilon_\lambda&B(\lambda;0,\varepsilon,1,(\lambda^2-1)/16)
   &\lambda\ne0,-1,\ \varepsilon=0,1\\
 G^\varepsilon&B(1;0,\varepsilon,0,1)&\varepsilon=0,1\\
 K&B(1;0,1,0,0)&
\end{array}
\]
\end{proposition}
\begin{proof}
$F^\varepsilon$ のAS4実現を仮定する。
三次・四次商座標の全係数の分母は
$\lambda,\lambda-1,\lambda+1,\lambda^2+1$ の積の冪だけであり、
全関係を消し自由列上でidentityである。
従ってその四因子を避けた各AS4値で、次元20,35の真の商基底となる。
中心三次行列の選んだ20次minorは
\[
 \frac{\lambda^6(\lambda-1)^{4+\varepsilon}(\lambda+1)^{6+\varepsilon}
 (\lambda^2+1)^{2-\varepsilon}(\lambda^2-\lambda+1)(\lambda^2+\lambda+1)}
 {2^{40+2\varepsilon}}.
\]
分母の極とこのminorの零点を併せた残余は、許容領域内では
\[
 \lambda=1,\qquad \lambda^2+1=0,\qquad
 \lambda^2+\lambda+1=0,\qquad\lambda^2-\lambda+1=0
\]
の七値である。
各 $\varepsilon$ について $1,i,\omega,-\omega$、
$\omega=(-1+\sqrt{-3})/2$ における全関係行列を別に構成し、
$\mathbb Q$, $\mathbb Q(i)$, $\mathbb Q(\sqrt{-3})$ 上で中心行列の階数20を確認する。
二次の三条件では係数共役がもう一方の根を覆う。
これにより商座標に極がある値も落とさず、全 $F^\varepsilon$ の中心三次は零となる。

三固定表示 $G^0,G^1,K$ についても、$\mathbb Q$ 上の中心三次行列の
非零20次minorを独立に確認した。
全表示で $R A R^{\mathsf T}$ の非零反巡回性を検算しているため、
AS4実現では $\mu=I$ である。従って中心三次の消滅がH3を排除する。
\end{proof}

Ore/PBW・normalizerと共鳴単射性は
\texttt{verify\_\allowbreak{}cubic\_\allowbreak{}as4\_\allowbreak{}four\_\allowbreak{}simple\_\allowbreak{}first.py}、
二つの全値rank coverと三固定表示は
\texttt{verify\_\allowbreak{}cubic\_\allowbreak{}as4\_\allowbreak{}four\_\allowbreak{}simple\_\allowbreak{}first\_\allowbreak{}remaining.py}
で検算する。この七表示の排除を、全共有非零領域の被覆と混同しない。
